\documentclass[10pt,a4paper]{amsart}
\usepackage{amsmath,amssymb,amsthm,graphicx}
% Retained minimal preamble for the Stacks Project chapter varieties.tex.
% Verbatim excerpt of the chapter preamble: Xy-pic setup (native preamble lines 14-18), the native theorem
% declarations with the mathematical macros (native preamble lines 154-217), and the project own no-op
% definitions of the reference/slogan/history display environments (native wrapper.tex lines 1-3) that the
% counted statement uses for its slogan. Copied mechanically; nothing added, removed or redefined.
% For commutative diagrams we use Xy-pic
\usepackage[all]{xy}

% We use 2cell for 2-commutative diagrams.
\xyoption{2cell}
% Theorem environments.
%
\theoremstyle{plain}
\newtheorem{theorem}[subsection]{Theorem}
\newtheorem{proposition}[subsection]{Proposition}
\newtheorem{lemma}[subsection]{Lemma}

\theoremstyle{definition}
\newtheorem{definition}[subsection]{Definition}
\newtheorem{example}[subsection]{Example}
\newtheorem{exercise}[subsection]{Exercise}
\newtheorem{situation}[subsection]{Situation}

\theoremstyle{remark}
\newtheorem{remark}[subsection]{Remark}
\newtheorem{remarks}[subsection]{Remarks}

\numberwithin{equation}{subsection}

% Macros
%
\def\lim{\mathop{\mathrm{lim}}\nolimits}
\def\colim{\mathop{\mathrm{colim}}\nolimits}
\def\Spec{\mathop{\mathrm{Spec}}}
\def\Hom{\mathop{\mathrm{Hom}}\nolimits}
\def\Ext{\mathop{\mathrm{Ext}}\nolimits}
\def\SheafHom{\mathop{\mathcal{H}\!\mathit{om}}\nolimits}
\def\SheafExt{\mathop{\mathcal{E}\!\mathit{xt}}\nolimits}
\def\Sch{\mathit{Sch}}
\def\Mor{\mathop{\mathrm{Mor}}\nolimits}
\def\Ob{\mathop{\mathrm{Ob}}\nolimits}
\def\Sh{\mathop{\mathit{Sh}}\nolimits}
\def\NL{\mathop{N\!L}\nolimits}
\def\CH{\mathop{\mathrm{CH}}\nolimits}
\def\proetale{{pro\text{-}\acute{e}tale}}
\def\etale{{\acute{e}tale}}
\def\QCoh{\mathit{QCoh}}
\def\Ker{\mathop{\mathrm{Ker}}}
\def\Im{\mathop{\mathrm{Im}}}
\def\Coker{\mathop{\mathrm{Coker}}}
\def\Coim{\mathop{\mathrm{Coim}}}

% Boxtimes
%
\DeclareMathSymbol{\boxtimes}{\mathbin}{AMSa}{"02}

%
% Macros for moduli stacks/spaces
%
\def\QCohstack{\mathcal{QC}\!\mathit{oh}}
\def\Cohstack{\mathcal{C}\!\mathit{oh}}
\def\Spacesstack{\mathcal{S}\!\mathit{paces}}
\def\Quotfunctor{\mathrm{Quot}}
\def\Hilbfunctor{\mathrm{Hilb}}
\def\Curvesstack{\mathcal{C}\!\mathit{urves}}
\def\Polarizedstack{\mathcal{P}\!\mathit{olarized}}
\def\Complexesstack{\mathcal{C}\!\mathit{omplexes}}
% \Pic is the operator that assigns to X its picard group, usage \Pic(X)
% \Picardstack_{X/B} denotes the Picard stack of X over B
% \Picardfunctor_{X/B} denotes the Picard functor of X over B
\def\Pic{\mathop{\mathrm{Pic}}\nolimits}
\def\Picardstack{\mathcal{P}\!\mathit{ic}}
\def\Picardfunctor{\mathrm{Pic}}
\def\Deformationcategory{\mathcal{D}\!\mathit{ef}}
\newenvironment{reference}{}{}
\newenvironment{slogan}{}{}
\newenvironment{history}{}{}
\usepackage{cleveref}
\crefname{originaldefinition}{original definition}{original definitions}
\Crefname{originaldefinition}{Original definition}{Original definitions}
\makeatletter
\newlabel{morphisms-definition-scheme-theoretic-image}{{01R7}{source}{}{sourcepointer-1}{}}
\newlabel{morphisms-definition-scheme-theoretic-image@cref}{{[originaldefinition][0][]01R7}{[1][0][]source}{}{}{}}
\makeatother
\crefname{originallemma}{original lemma}{original lemmas}
\Crefname{originallemma}{Original lemma}{Original lemmas}
\makeatletter
\newlabel{morphisms-lemma-quasi-compact-scheme-theoretic-image}{{01R8}{source}{}{sourcepointer-2}{}}
\newlabel{morphisms-lemma-quasi-compact-scheme-theoretic-image@cref}{{[originallemma][0][]01R8}{[1][0][]source}{}{}{}}
\makeatother
\title{The scheme theoretic image of a product of quasi-compact morphisms is the product of the scheme theoretic images}
\author{The Stacks Project Authors}
\date{stacks-project:varieties.tex:local-snapshot-20261011; GNU FDL 1.2 or later}
\begin{document}
\maketitle
\noindent\textit{Editorial note.} All delivered mathematical text is the author's own, copied from the pinned Stacks Project chapter file varieties.tex; no statement or proof step was rewritten, reordered or removed.\par\medskip
\noindent\textit{Editorial note.} The proof ends with the author's own 'Details omitted.' after the displayed tensor product equality; the sentence is reproduced verbatim. The omitted verification is the bounded bookkeeping comparing the two kernels, and the reduction that isolates it is complete as written.\par\medskip
\noindent\textit{Editorial note.} The counted spans carry 2 navigation stubs and no bibliography citation.\par\medskip
\noindent\textit{Editorial note.} The preamble is the chapter's retained minimal preamble (sources/stacks-varieties/retained-preamble.tex lines 1-77), a verbatim excerpt of the native preamble.\par\medskip
\noindent\textit{Editorial note.} The label lemma-scheme-theoretic-image-product-map is defined only in varieties.tex, so no cross-chapter restatement had to be weighed for this card.\par\medskip
\noindent\textit{Allowed standard prerequisites (editorial).}\begin{itemize}
\item Scheme theoretic image as the smallest closed subscheme through which the morphism factors. Morphisms, Definition morphisms-definition-scheme-theoretic-image, invoked at proof lines 4491-4492 and named in the statement. Supplies the universal property used for the first inclusion and the kernel description in the affine case.
\item Scheme theoretic images are local on affine opens and computed as kernels of restriction maps. Morphisms, Lemma morphisms-lemma-quasi-compact-scheme-theoretic-image, invoked at proof lines 4498-4502. Supplies the reduction to affine X and Y and the kernel description of the images.
\item Global functions on a product of affine schemes over a field. Standard, used at proof lines 4512-4516. Supplies the equality of the global sections of A times\_k B with the tensor product of the coordinate rings, which yields the reverse inclusion.
\end{itemize}\medskip
\ifdefined\hypertarget\hypertarget{sourcepointer-1}{}\fi
\noindent\textit{Source reference.} Original definition 01R7: sources/stacks-chow/tags.txt, source lines 2292--2292. Pointer outside the counted span, invoked once for context; the load-bearing argument of the counted statement is the author's proof reproduced in full in the proof span.\par\medskip
\ifdefined\hypertarget\hypertarget{sourcepointer-2}{}\fi
\noindent\textit{Source reference.} Original lemma 01R8: sources/stacks-chow/tags.txt, source lines 2293--2293. Pointer outside the counted span, invoked once for context; the load-bearing argument of the counted statement is the author's proof reproduced in full in the proof span.\par\medskip
\begin{lemma}
\label{lemma-scheme-theoretic-image-product-map}
Let $k$ be a field.
Let $f : A \to X$, $g : B \to Y$ be quasi-compact morphisms of schemes
over $k$. Let $Z \subset X$ be the scheme theoretic image of $f$, see
Morphisms, Definition \ref{morphisms-definition-scheme-theoretic-image}.
Similarly, let $Z' \subset Y$ be the scheme theoretic image of $g$.
Then $Z \times_k Z'$ is the scheme theoretic image of $f \times g$.
\end{lemma}

\begin{proof}
Recall that $Z$ is the smallest closed subscheme of $X$ through which
$f$ factors. Similarly for $Z'$. Let $W \subset X \times_k Y$ be the
scheme theoretic image of $f \times g$. As $f \times g$ factors through
$Z \times_k Z'$ we see that $W \subset Z \times_k Z'$.

\medskip\noindent
To prove the other inclusion let $U \subset X$ and $V \subset Y$ be
affine opens. By
Morphisms, Lemma \ref{morphisms-lemma-quasi-compact-scheme-theoretic-image}
the scheme $Z \cap U$ is the scheme theoretic image of
$f|_{f^{-1}(U)} : f^{-1}(U) \to U$, and similarly for
$Z' \cap V$ and $W \cap U \times_k V$. Hence we may assume
$X$ and $Y$ affine. As $f$ and $g$ are quasi-compact this implies
that $A = \bigcup U_i$ is a finite union of affines and
$B = \bigcup V_j$ is a finite union of affines.
Then we may replace $A$ by $\coprod U_i$ and $B$ by
$\coprod V_j$, i.e., we may assume that $A$ and $B$ are affine as well.
In this case $Z$ is cut out by
$\Ker(\Gamma(X, \mathcal{O}_X) \to \Gamma(A, \mathcal{O}_A))$
and similarly for $Z'$ and $W$. Hence the result follows from
the equality
$$
\Gamma(A \times_k B, \mathcal{O}_{A \times_k B})
=
\Gamma(A, \mathcal{O}_A) \otimes_k \Gamma(B, \mathcal{O}_B)
$$
which holds as $A$ and $B$ are affine. Details omitted.
\end{proof}

\end{document}
